\scanpage{7}
% Lines 1–2. The faint reversed typescript shows through from the other face.
formellement sur le \inserted{\unsure{un|$\sim$}} diagramme de deux suites transposées l'une de l'autre
% Lines 3–4
\diagram{\begin{tikzcd}[ampersand replacement=\&,column sep=large]
 \unsure{A|A_1} \arrow[r,"F"] \& A^{(p)}
 \arrow[r,phantom,"{\overset{\unsure{\mathbb{V}|V}}{\unsure{\relbar|\to}}}" description] \& V\\
 B \& B^{(p)} \arrow[l,"V"'] \& B \arrow[l,"F"']
\end{tikzcd}}

% Lines 5–7
On \unsure{trouve|a donc} ainsi la suite exacte
\diagram{\begin{tikzcd}[ampersand replacement=\&,column sep=large]
 \unsure{0|\cdot} \arrow[r,dash] \& \operatorname{Gr}(\underline{t}_A) \arrow[r,-{Stealth[length=2pt,width=2pt]}] \& {}_p A \arrow[r,-{Stealth[length=2pt,width=2pt]}] \& D(\operatorname{Gr}(\underline{t}_B)) \arrow[r] \& 0
\end{tikzcd}}
extension canonique de $D({}_F\deleted{\unsure{A}}B)$ par $ {}_F A$.

% Lines 8–12. The enclosed interline has a point descending after “abéliennes”.
Noter d'ailleurs que $S_A$ est une extension de $p$-alg. de Lie abéliennes \inserted{$t_A^\vee$ \unsure{étant|est} muni de la puiss. \unsure{$p$.ième|$p$-ième} nulle.}, et correspond donc en fait à une extension de groupes \unsure{radiciels}. Cette extension peut-elle se déduire de la connaissance de $ {}_p A$ ??
% archive sheet number: 6
